\documentclass[12pt,a4paper]{article}

\usepackage[margin=1in]{geometry}
\usepackage[T1]{fontenc}
\usepackage[utf8]{inputenc}
\usepackage{setspace}
\usepackage{natbib}
\usepackage[hidelinks]{hyperref}
\usepackage{csquotes}
\usepackage{amsmath}
\usepackage{amssymb}
\usepackage{booktabs}
\usepackage{graphicx}
\usepackage{microtype}
\microtypesetup{expansion=false}
\usepackage{float}

\setstretch{1.15}
\raggedbottom
\bibliographystyle{apalike}
\graphicspath{
  {../../outputs/fixation_circular_working_memory/figures/}
  {../../outputs/fixation_circular_working_memory/perturbation_experiments/noise_structure/gaussian_control/figures/}
}

\title{Recovering Psilocybin-Related Working-Memory Signatures by Perturbing a Recurrent Neural Network}
\author{}
\date{}

\begin{document}
\maketitle

\begin{abstract}
Acute psilocybin produces heterogeneous effects across human working-memory and
attention tasks, including response slowing, greater impairment at higher
updating load, and vulnerability to competing information. This dissertation
tests whether such a selective behavioural signature can be recovered by
perturbing a recurrent neural network (RNN) trained to perform working-memory
tasks. We first established a fixation-gated continuous-time RNN that
maintained a circular stimulus through variable delays, replicated its
performance across five independently trained models, and evaluated
independent Gaussian hidden-state noise as a non-specific degradation control.
Mean response error was
$4.36\pm1.55^{\circ}$ and fixation accuracy was $0.9677\pm0.0016$ across
models; hidden-state decoding recovered delay-period angle with
$0.55\pm0.46^{\circ}$ error. Gaussian noise caused graded, delay-dependent
impairment while largely preserving fixation and circular centroid spacing.
These results establish a reproducible baseline and quantify generic
stochastic disruption. Gaussian noise caused broad, delay-dependent
degradation rather than recovering a selective human behavioural signature. It
therefore provides a control for evaluating literature-motivated perturbations,
without being interpreted as a biological model of psilocybin.
\end{abstract}

\section{Introduction}

Working memory depends on retaining task-relevant information after its
eliciting input has disappeared. Recurrent dynamics provide a natural
computational substrate for this function: activity can settle onto an
attractor or evolve along a slow manifold that preserves a decodable memory
variable \citep{wang2013nmda,kilpatrick2013,ghazizadeh2021}. For a continuous
circular variable, the relevant question is not only whether a network produces
an accurate response, but whether its hidden population maintains ordered
geometry, resists drift, and generalises beyond the delays encountered during
training.

The dissertation's central question is whether perturbing this recurrent
system can reproduce the qualitative pattern of working-memory effects
observed in humans during acute psilocybin administration. Human findings are
mixed rather than uniformly impairing. Psilocybin has produced greater
impairment under higher N-back updating load \citep{barrett2018}, while a small
study found clearer impairment in multiple-object tracking than in spatial
working-memory span \citep{carter2005}. The behavioural target is therefore a
selective signature---including slowing, load sensitivity, and susceptibility
to competing information---rather than an arbitrary reduction in accuracy.

Psychedelic neuroscience motivates candidate dynamical explanations for this
signature. The
Entropic Brain hypothesis proposes that psychedelic states occupy a broader,
less constrained repertoire of brain states \citep{carhartharris2014entropic}.
Human MEG, EEG, and fMRI studies have reported increased signal diversity or
altered network entropy under psilocybin and related compounds
\citep{schartner2017increased,timmermann2019neural,viol2017shannon}. These
measurements are direct evidence about recorded brain signals, but they do not
imply that psilocybin is equivalent to additive noise in a task-trained RNN.
Increased diversity could reflect structured changes in responsiveness,
coupling, context dependence, or state-space exploration.

REBUS (``relaxed beliefs under psychedelics'') places this change within
predictive processing \citep{carhartharris2019rebus}. The theory proposes
reduced precision of high-level priors and greater influence of bottom-up
signals, relating precision to postsynaptic gain. A working-memory translation
is plausible because a maintained representation must resist irrelevant input.
However, scaling an RNN's input drive, recurrent drive, or activation slope is
an implementation hypothesis; REBUS does not prescribe any of these equations.
The distinction matters because the three manipulations can make different
predictions about clean maintenance, distractor capture, and recovery.

A more direct computational precedent comes from the receptor-informed
whole-brain model of \citet{herzog2023mechanistic}. That study represented
5-HT2A receptor activation as regionally heterogeneous modulation of the
frequency--current response function. Response gain interacted with structural
connectivity to produce non-uniform entropy changes. The study therefore
supports response-gain modulation as a candidate modelling dimension and
suggests that recurrent topology shapes its consequences. It did not study
working memory, a trained RNN, or a pharmacological mapping for the abstract
gain values used here.

The direction and population specificity of gain changes cannot be inferred
from that whole-brain result alone. In mouse visual cortex, the hallucinogenic
5-HT2A agonist DOI reduced visually evoked response gain
\citep{michaiel2019visualgain}; cell-type-specific manipulation of a
5-HT2A-related pathway produced heterogeneous effects and divisive suppression
through balanced excitation and inhibition \citep{barzan2024gaincontrol}.
Separately, human EEG modelling linked perceptual precision to the synaptic
gain of superficial pyramidal prediction-error units
\citep{brown2012precision}. These studies support bidirectional,
pathway-specific tests, not a presumption that a psychedelic-like manipulation
is a uniform gain increase.

We hypothesise that a literature-motivated perturbation of a trained
working-memory RNN can recover the qualitative human signature more closely
than generic stochastic disruption. This requires more than decreasing
performance: a successful perturbation should reproduce a selective pattern
across clean maintenance, delay duration, response dynamics, updating demand,
and competing input. The present analysis establishes the unperturbed baseline
and tests independent Gaussian noise as the generic control. Gain modulation
is retained as a candidate mechanism because it can alter response sensitivity
or the balance of external and recurrent drive, but any such manipulation is
an abstract computational hypothesis rather than a biological equivalent of
psilocybin.

\section{Methods}

\subsection{Fixation-gated circular task}

Each stimulus angle $\theta\in[0,2\pi)$ was encoded over 32 evenly spaced
preferred angles $\mu_j$ using
\begin{equation}
  c_j(\theta)=\exp\{\kappa[\cos(\theta-\mu_j)-1]\},\qquad \kappa=8.
\end{equation}
Trials comprised pre-cue fixation, cue, delay, and response periods. Pre-cue
duration was sampled from 15, 25, or 35 steps; cue duration from 10, 20, or 30
steps; and delay duration from 10, 20, 40, or 80 steps. The response lasted 25
steps. A fixation input was one until response onset and zero thereafter. The
target fixation output followed the same schedule; the 32 circular outputs
were silent before response and reported the remembered angle during response.
Thus delay-period memory was evaluated from hidden state rather than from the
deliberately silent circular readout.

\subsection{Network and implemented recurrence}

The model contained 33 inputs, 64 recurrent units, and 33 outputs. With
$\alpha=dt/\tau=20/100=0.2$, the code implemented
\begin{align}
  z_t &= W_{\mathrm{in}}x_t+b_{\mathrm{in}}
       +W_{\mathrm{rec}}h_{t-1}+b_{\mathrm{rec}},\\
  h_t &= \tanh\!\left[(1-\alpha)h_{t-1}+\alpha z_t+\varepsilon_t\right],\\
  y_t &= W_{\mathrm{out}}h_t+b_{\mathrm{out}} .
\end{align}
This placement of the leak, perturbation, and non-linearity is the implemented
update; it should not be replaced by a textbook continuous-time equation.
During training only, independent recurrent noise with standard deviation
0.05 was added after the leaky update and before $\tanh$.

\subsection{Training and replication}

We trained with Adam (learning rate $10^{-3}$) for 4,000 steps using batches of
64 trials. Gaussian input noise had standard deviation 0.01. The weighted
mean-squared error assigned weight 5 to circular response output and weight 2
to fixation output. The first five trial steps and the first five response
transition steps were unscored. Gradients were clipped to a maximum norm/value
of 1.0 as implemented by the training utility. Five networks were trained from
independent seeds 20260714--20260718. Model seed, rather than trial or
stochastic batch, was the replication unit.

\subsection{Baseline evaluation and dynamical analyses}

Each network was evaluated over 20 batches. Response angle was decoded by
population-vector averaging, and absolute circular error was the shortest
wrapped angular distance from target. Fixation accuracy measured correct
wait--go output state. Generalisation sweeps used delays of 10, 20, 40, 80, and
160 steps.

Delay-period content was decoded from hidden states by ridge regression to the
two-dimensional target vector
\[
  \bigl[\cos(\theta),\,\sin(\theta)\bigr].
\]
The decoder used 512 trials, a 75:25
train--test split, and ridge coefficient 1.0. Cross-temporal matrices fitted a
decoder at each training time and tested it at each held-out time.

For the canonical checkpoint only, late-delay states seeded blank-input
autonomous probes and fixed-point optimisation. Local Jacobians were computed
at the resulting approximate fixed points. We report residual, decoded angle,
drift, eigenvalue magnitude, and alignment of the leading eigenvector with the
sampled circular tangent. These sampled analyses are descriptive properties of
one checkpoint, not five-model estimates or proof of a global attractor.

\subsection{Independent Gaussian control}

Frozen networks received perturbations only during the delay. For each active
time $t$, trial $b$, and hidden unit $i$, an independent
$\eta_{tbi}\sim\mathcal{N}(0,1)$ was drawn and normalised as
\begin{equation}
\varepsilon_{tbi}=
\sigma\frac{\eta_{tbi}}
{\left(|\mathcal{A}|^{-1}\sum_{(t,b,i)\in\mathcal{A}}\eta_{tbi}^{2}\right)^{1/2}},
\label{eq:gaussian}
\end{equation}
where $\mathcal{A}$ contains all delay-active entries. The realised within-batch
RMS therefore equalled target strength $\sigma$. We evaluated delays 20, 80,
and 160 at
\[
  \sigma\in\{0,\ 0.01,\ 0.025,\ 0.05,\ 0.10\}.
\]
Each model--delay condition
used five stochastic replicates of 20 paired 64-trial batches. Clean-fitted
decoders quantified representational distortion without adapting to the noisy
states. RMS is an intervention magnitude, not a pharmacological dose.

Outcomes were response error, clean-decoder error, angular drift, fixation
accuracy, circular-state geometry, and response settling. Geometry was
measured in a clean-fitted principal-component plane using within-angle
dispersion and between-angle centroid spacing. Settling was the first scored
response step with error below $20^{\circ}$ after excluding the five transition
steps.

\section{Results}

\subsection{The baseline replicated accurate circular working memory}

Across five independently trained models, mean response angular error was
$4.36\pm1.55^{\circ}$ (range $2.95$--$7.03^{\circ}$), response population MSE
was $0.00957\pm0.00194$, and fixation accuracy was
$0.9677\pm0.0016$ (mean $\pm$ SD across model seeds). Delay-period diagonal
decoding error was $0.55\pm0.46^{\circ}$ (range
$0.31$--$1.37^{\circ}$), showing that the remembered angle remained linearly
recoverable despite the silent circular output.

\begin{figure}[p]
\centering
\includegraphics[width=.92\linewidth,height=.68\textheight,keepaspectratio]{fixation_circular_working_memory_seed_sweep_delay_curves.png}
\caption{\textbf{Baseline delay generalisation across five independently
trained networks.} Response angular error is shown for trained delays
10--80 steps and the untrained 160-step delay. Curves identify model seeds,
which are the replication unit. Performance remains substantially above a
random circular response at long delays, while between-seed variation increases.}
\label{fig:baseline-delay}
\end{figure}

The five models generalised to the untrained 160-step delay, with seed-level
mean response errors ranging from $5.14^{\circ}$ to $11.73^{\circ}$
(Fig.~\ref{fig:baseline-delay}). This increase indicates accumulating
delay-dependent error rather than perfect time invariance.

\begin{figure}[p]
\centering
\includegraphics[width=.90\linewidth,height=.68\textheight,keepaspectratio]{fixation_circular_working_memory_cross_temporal_decoder.png}
\caption{\textbf{Cross-temporal hidden-state decoding for the canonical
checkpoint.} Ridge decoders predict sine and cosine of target angle from held-out
hidden states. The low delay-period diagonal error
($0.427^{\circ}$; maximum $0.474^{\circ}$) confirms precise maintenance; the
full matrix describes how the representational code transfers across time.}
\label{fig:baseline-decoder}
\end{figure}

\subsection{A canonical checkpoint showed sampled attractor-like geometry}

For the canonical checkpoint, fixed-point optimisation produced a mean
one-step residual of 0.00202, with 28.1\% of sampled endpoints below the
pre-specified 0.001 residual threshold. Mean decoded fixed-point error was
$1.24^{\circ}$ and mean drift from the corresponding late-delay state was
$1.10^{\circ}$. The mean leading eigenvalue magnitude was 0.99975, the mean
second-largest magnitude was 0.88742, and leading-eigenvector alignment with
the sampled circular tangent was 0.978.

\begin{figure}[p]
\centering
\includegraphics[width=.88\linewidth,height=.68\textheight,keepaspectratio]{fixation_circular_working_memory_fixed_point_analysis.png}
\caption{\textbf{Sampled fixed-point and Jacobian analysis of one canonical
checkpoint.} Optimisation residuals, hidden-decoder angle preservation, drift,
and local eigenvalue magnitudes are shown for late-delay starts. A near-neutral
leading direction with contracting secondary dynamics is consistent with a
ring-like slow manifold. The finite sample and incomplete convergence preclude
an exhaustive claim about the full state space.}
\label{fig:baseline-dynamics}
\end{figure}

Together, preserved angle, low drift, and a near-neutral tangent direction
support an attractor-like or slow-manifold account
(Fig.~\ref{fig:baseline-dynamics}). They remain descriptive because the
analysis used one checkpoint and sampled initial states.

\subsection{Gaussian noise caused weak, delay-dependent degradation}

Gaussian noise produced graded impairment (Fig.~\ref{fig:gaussian-behaviour}).
At RMS 0.05, response error increased from
$4.34\pm1.54^{\circ}$ to $4.91\pm1.45^{\circ}$ at delay 20, from
$6.49\pm1.93^{\circ}$ to $7.47\pm1.68^{\circ}$ at delay 80, and from
$9.03\pm2.50^{\circ}$ to $10.44\pm2.24^{\circ}$ at delay 160. At RMS 0.10,
the corresponding errors were $6.42\pm1.35^{\circ}$,
$10.48\pm1.92^{\circ}$, and $14.68\pm2.57^{\circ}$.

\begin{figure}[p]
\centering
\includegraphics[width=.94\linewidth,height=.67\textheight,keepaspectratio]{figure_g1_dose_delay_behaviour.pdf}
\caption{\textbf{Response error under independent Gaussian delay
perturbation.} Absolute circular error is shown across RMS strengths and delay
lengths. Solid lines are five-model means; dashed lines and shading are 95\%
bootstrap intervals across models. RMS quantifies perturbation magnitude, not
drug dose.}
\label{fig:gaussian-behaviour}
\end{figure}

Maintenance measures were more sensitive than final response. At delay 80 and
RMS 0.05, clean-decoder error increased from
$1.29\pm0.82^{\circ}$ to $3.71\pm1.53^{\circ}$ and decoded drift from
$2.60\pm1.84^{\circ}$ to $5.20\pm2.11^{\circ}$.
Fixation accuracy remained $0.98\pm0.00$ in both conditions
(Fig.~\ref{fig:gaussian-maintenance}).

\begin{figure}[p]
\centering
\includegraphics[width=.94\linewidth,height=.67\textheight,keepaspectratio]{figure_g2_maintenance_metrics.pdf}
\caption{\textbf{Maintenance outcomes at delay 80.} Clean-fitted decoder error
and decoded drift increase with perturbation RMS, whereas fixation accuracy is
preserved. Lines and intervals are defined as in
Fig.~\ref{fig:gaussian-behaviour}.}
\label{fig:gaussian-maintenance}
\end{figure}

At delay 80 and RMS 0.05, within-angle dispersion increased from 0.550 to
0.732, whereas normalised between-angle centroid spacing remained 0.990
relative to the clean value of one (Fig.~\ref{fig:gaussian-geometry}). The
representation therefore blurred without a wholesale collapse of its circular
organisation.

\begin{figure}[p]
\centering
\includegraphics[width=.92\linewidth,height=.67\textheight,keepaspectratio]{figure_g3_geometry.pdf}
\caption{\textbf{Clean and Gaussian circular-state geometry at delay 80 and
RMS 0.05.} Angle centroids and uncertainty ellipses are projected into the
clean-fitted PC1--PC2 plane. Gaussian perturbation broadens within-angle states
while largely preserving global centroid separation.}
\label{fig:gaussian-geometry}
\end{figure}

Response settling changed little at the primary RMS. At delay 80, the first
response step below $20^{\circ}$ increased from 5.02 clean to 5.19 at RMS 0.05,
with 99.2\% of trials settling. At RMS 0.10, settling increased to 6.55 steps
and 93.6\% settled (Fig.~\ref{fig:gaussian-settling}). Final error in this
focused analysis was $5.94^{\circ}$ clean, $6.75^{\circ}$ at RMS 0.05, and
$10.40^{\circ}$ at RMS 0.10.

\begin{figure}[p]
\centering
\includegraphics[width=.92\linewidth,height=.67\textheight,keepaspectratio]{figure_g4_settling.pdf}
\caption{\textbf{Response settling at delay 80.} Settling is the first counted
response step below $20^{\circ}$ error after excluding the five transition
steps. Points are model seeds and lines are five-model means.}
\label{fig:gaussian-settling}
\end{figure}

Independent Gaussian noise is therefore a useful weak, non-specific
degradation control. It increases dispersion, drift, and error with elapsed
delay, but does not identify a pharmacological or cognitive mechanism.
Structured AR(1) and topology-correlated noise remain a separate exploratory
robustness study and are not part of the present mechanism comparison.

\section{Discussion}

The fixation-gated RNN provides a replicated continuous working-memory
baseline. Accurate response and delay-period decoding across five seeds show
that the result is not restricted to one favourable initialisation. The
canonical checkpoint's sampled fixed points and local Jacobians are consistent
with a ring-like slow manifold, although that inference is descriptive and
does not substitute for replicated dynamical analysis.

The dissertation asks whether perturbing a working-memory RNN can recover the
selective behavioural signature observed in humans under acute psilocybin.
The present results do not yet answer that hypothesis, because they establish
the baseline and generic control rather than a successful
psilocybin-motivated manipulation. They nevertheless provide the necessary
reference against which recovery of the human signature can be judged.

Independent Gaussian perturbation supplied the intended generic control. Its
effects were graded and strongest in maintenance-sensitive measures, while
fixation and between-angle organisation were largely preserved at moderate
strength. This establishes how unstructured variability degrades the trained
system. It does not establish that altered entropy in psychedelic recordings
is Gaussian, nor that RMS maps to human dose. More importantly, the principal
pattern was a broad increase in error with perturbation strength and delay.
This is insufficient by itself to recover the task-selective combination of
slowing, updating-load sensitivity, and competing-input vulnerability that
defines the dissertation hypothesis.

The gain literature supplies one possible mechanistic interpretation, but not
a completed result. Receptor-informed modelling supports altered neural
response gain \citep{herzog2023mechanistic}, whereas REBUS motivates altered
precision and effective afferent influence
\citep{carhartharris2019rebus}. Direct sensory-cortex studies also caution that
gain effects can be reduced, heterogeneous, and population specific
\citep{michaiel2019visualgain,barzan2024gaincontrol}. Consequently, any later
RNN gain manipulation must be judged by whether it recovers the human
behavioural pattern, not by whether it simply changes network activity or
increases error.

\section{Conclusion}

The fixation-gated RNN provides a reproducible substrate for testing the
dissertation's central hypothesis. Independent Gaussian noise produces
predictable delay-dependent degradation, but the current evidence does not
show that it recovers the selective behavioural signatures reported in humans
under acute psilocybin. The relevant standard for subsequent perturbations is
therefore pattern recovery across working-memory conditions, not impairment
alone. Even if an RNN perturbation meets that standard, it would demonstrate
computational sufficiency rather than biological equivalence to psilocybin.

\bibliography{references}

\end{document}
